\chapter{System Design}
\label{ch:design}

\section{Introduction}
The design follows the boundaries established in Chapter~\ref{ch:analysis}:
visitor search, an inspectable selected query, an independent managed catalogue,
and a local administrator. The diagrams in this chapter describe the implemented
application. They do not introduce a multi-store marketplace, cloud inference,
or an identity model that is absent from the code.

\section{Architectural Design}
\subsection{Client, API, and Model Services}
Figure~\ref{fig:architecture} shows the browser, FastAPI service, neural services,
and data stores. The browser handles file selection, canvas interaction, preview,
and result presentation. Server-side Python validates uploads, performs neural
inference, ranks vectors, and mediates database changes. FastAPI serves both
the API and local frontend assets under one origin.

\begin{figure}[H]\centering
\begin{tikzpicture}[node distance=7mm]
\node[reportbox,text width=12.4cm] (browser) {Browser: image input, selection canvas, query preview, results, product management};
\node[reportbox,text width=12.4cm,below=of browser] (api) {FastAPI: request validation and route handlers\\Local administrator session, origin and CSRF checks};
\node[reportbox,text width=5.5cm,minimum height=20mm,below left=8mm and -5.7cm of api] (sam) {SAM2.1 segmentation\\Points to mask and overlay\\White-background cutout};
\node[reportbox,text width=5.5cm,minimum height=20mm,below right=8mm and -5.7cm of api] (encode) {SigLIP2 image encoding\\Shared preprocessing\\Normalised vector ($d=768$)};
\node[storebox,text width=5.5cm,minimum height=23mm,below=10mm of sam] (gallery) {Original research gallery\\3,068 images and metadata\\Persistent FAISS IndexFlatIP};
\node[storebox,text width=5.5cm,minimum height=23mm,below=10mm of encode] (managed) {Added-product catalogue\\SQLite products, photos, vectors\\Per-request FAISS ranking\\One result per product};
\node[storebox,text width=12.4cm,below=10mm of $(gallery.south)!0.5!(managed.south)$] (local) {Independent local stores\\Administrator sessions; search telemetry; backup archives};
\draw[reportarrow] (browser)--(api);
\draw[reportarrow] (api.south -| sam.north)--(sam.north);
\draw[reportarrow] (api.south -| encode.north)--(encode.north);
\draw[reportarrow] (sam.east)--(encode.west);
\draw[reportarrow] (encode.south)--(managed.north);
\draw[reportarrow] (encode.south west)--(gallery.north east);
\draw[reportarrow,dashed] (api.east)--++(0.45,0) |- (local.east);
\end{tikzpicture}

\caption{Implemented local architecture and separate catalogue resources}
\label{fig:architecture}\end{figure}

The application loads the checkpoints locally at startup. The health endpoint
reports readiness, devices, active catalogue source, and indexed image count.
The launcher waits for a compatible healthy service before opening the browser
and reuses an already-ready instance instead of starting another server on the
same port. The service binds to the loopback interface for the current deployment.

\subsection{Separation of the Two Catalogues}
The original gallery uses image files, a versioned metadata mapping, and a
persistent FAISS index. The added-product catalogue uses SQLite records containing
product attributes, encoded photographs, and vectors. Updating a managed product
does not rewrite the research dataset, its identity review, or the persistent
gallery index. This separation protects the evidence base while permitting
practical product entry.

The managed search builds an in-memory exact FAISS index from available product
images for each request. Its filtering and product aggregation therefore differ
from the original-gallery path. The common elements are the fixed image encoder,
normalisation, query preparation, and user selection tools.

\section{Behavioural Design}
\subsection{Use Case Diagram}
Figure~\ref{fig:usecases} groups visitor and administrator interactions. The
administrator may also use visitor search; the diagram distinguishes privileged
operations rather than claiming that an administrator is prevented from searching.
All management routes are checked server-side even when called without the
management page.

\begin{figure}[H]\centering\begin{tikzpicture}[x=1cm,y=1cm]
\node[reportbox,text width=2cm] (visitor) at (-5.1,2.8) {Visitor};
\node[reportbox,text width=2cm] (admin) at (5.1,-2) {Local administrator};
\node[reportbox,ellipse,text width=5.1cm] (search) at (0,3.5) {Search either catalogue};
\node[reportbox,ellipse,text width=5.1cm] (select) at (0,1.8) {Select / refine / preview a garment};
\node[reportbox,ellipse,text width=5.1cm] (feedback) at (0,0.1) {Inspect results and provide feedback};
\node[reportbox,ellipse,text width=5.1cm] (manage) at (0,-1.6) {Create, edit, and delete products};
\node[reportbox,ellipse,text width=5.1cm] (backup) at (0,-3.3) {Export, preview, restore, and recover};
\node[reportbox,ellipse,text width=5.1cm] (account) at (0,-5) {Manage local account and reindex};
\node[draw=UniversityBlue!50,dashed,fit=(search)(select)(feedback)(manage)(backup)(account),inner sep=5mm,
 label={[font=\small]above:Local visual-search application}] {};
\draw (visitor)--(search);\draw (visitor)--(select);\draw (visitor)--(feedback);
\draw (admin)--(manage);\draw (admin)--(backup);\draw (admin)--(account);
\end{tikzpicture}

\caption{Visitor and local-administrator use cases}
\label{fig:usecases}\end{figure}

\subsection{Search Activity}
Figure~\ref{fig:search-activity} follows the accepted image from upload to result
review. Full-image mode bypasses segmentation. SAM2 and manual modes prepare
an alternative image; they do not switch to a different retrieval model. A
category correction or a revised selection creates a new query rather than
changing the meaning of an earlier result.

\begin{figure}[H]\centering\begin{tikzpicture}[node distance=5mm]
\node[reportbox,text width=8.1cm] (upload) {Upload and validate image; choose catalogue};
\node[reportbox,text width=8.1cm,below=of upload] (mode) {Use the full image or prepare a selected query};
\node[reportbox,text width=3.9cm,minimum height=22mm,below left=8mm and -3.5cm of mode] (full) {Full-image query};
\node[reportbox,text width=3.9cm,minimum height=22mm,below right=8mm and -3.5cm of mode] (selected) {SAM2 points or\\manual region\\Inspect and accept};
\node[reportbox,text width=8.1cm,below=8mm of $(full.south)!0.5!(selected.south)$] (preview) {Preview the exact submitted image};
\node[reportbox,text width=8.1cm,below=of preview] (encode) {Shared preparation; encode and normalise};
\node[reportbox,text width=8.1cm,below=of encode] (search) {Apply category rule; FAISS similarity ranking};
\node[reportbox,text width=8.1cm,below=of search] (display) {Group products where configured; show results};
\node[reportbox,text width=8.1cm,below=of display] (review) {Inspect matches; correct category, selection, or source};
\draw[reportarrow] (upload)--(mode);
\draw[reportarrow] (mode.south)--(full.north);\draw[reportarrow] (mode.south)--(selected.north);
\draw[reportarrow] (full.south)--(preview.north);\draw[reportarrow] (selected.south)--(preview.north);
\draw[reportarrow] (preview)--(encode);\draw[reportarrow] (encode)--(search);
\draw[reportarrow] (search)--(display);\draw[reportarrow] (display)--(review);
\end{tikzpicture}

\caption{Search activity with full-image and selected-query branches}
\label{fig:search-activity}\end{figure}

\subsection{Segmentation and Search Sequence}
The interaction in Figure~\ref{fig:sequence} has two explicit stages. A segment
request sends the original image and the accumulated point labels, and receives
the overlay and cutout. The user then accepts a query and submits a separate
search request. Thus a segmentation response is not itself a retrieval result.
Further points trigger another segmentation request, while resetting or replacing
the image invalidates earlier pending responses in the frontend.

\begin{figure}[H]\centering\begin{tikzpicture}[x=1cm,y=1cm]
\foreach \x/\name in {0/Browser,4/FastAPI,8/Model services,12/Gallery index} {
 \node[reportbox,text width=2.25cm,font=\small] at (\x,0) {\name};
 \draw[dashed,black!45] (\x,-0.5)--(\x,-9);
}
\draw[reportarrow] (0,-1)--node[above,font=\scriptsize]{image + point labels}(4,-1);
\draw[reportarrow] (4,-1.9)--node[above,font=\scriptsize]{validated SAM2 input}(8,-1.9);
\draw[reportarrow] (8,-2.8)--node[above,font=\scriptsize]{mask, overlay, cutout}(4,-2.8);
\draw[reportarrow] (4,-3.6)--node[above,font=\scriptsize]{selection preview}(0,-3.6);
\node[reportnote,text width=3cm] at (0,-4.3) {User refines or accepts};
\draw[reportarrow] (0,-5)--node[above,font=\scriptsize]{accepted query + category}(4,-5);
\draw[reportarrow] (4,-5.8)--node[above,font=\scriptsize]{preprocess + SigLIP2}(8,-5.8);
\draw[reportarrow] (8,-6.6)--node[above,font=\scriptsize]{normalised vector}(4,-6.6);
\draw[reportarrow] (4,-7.3)--node[above,font=\scriptsize]{exact ranking / category filter}(12,-7.3);
\draw[reportarrow] (12,-8.1)--node[above,font=\scriptsize]{ranked images and metadata}(4,-8.1);
\draw[reportarrow] (4,-8.9)--node[above,font=\scriptsize]{results + query hash}(0,-8.9);
\end{tikzpicture}

\caption{Request sequence for SAM2 selection followed by original-gallery search}
\label{fig:sequence}\end{figure}

\section{Component Design}
Figure~\ref{fig:components} identifies the principal Python classes. The
\texttt{SearchEngine} owns image encoding and the original index.
\texttt{Segmenter} owns point-prompted mask inference. \texttt{DemoCatalog}
implements the independent product collection. \texttt{CatalogBackups} validates
and restores its portable archives. \texttt{AdminAuth} manages one local account
and its sessions, while \texttt{SearchObservability} records search events and
feedback in a separate database.

\begin{figure}[H]\centering\begin{tikzpicture}[node distance=11mm]
\node[reportbox,text width=5.6cm,align=left] (engine) {
 \textbf{SearchEngine}\\\rule{5.3cm}{0.3pt}\\
 index; image paths; catalogue vectors\\
 \rule{5.3cm}{0.3pt}\\\texttt{encode\_image}, \texttt{search}\\
 \texttt{load\_index}, \texttt{build\_index}};
\node[reportbox,text width=5.6cm,align=left,right=of engine] (segmenter) {
 \textbf{Segmenter}\\\rule{5.3cm}{0.3pt}\\
 SAM2 processor; model; device\\
 \rule{5.3cm}{0.3pt}\\\texttt{segment}\\
 mask selection and cutout encoding};
\node[reportbox,text width=5.6cm,align=left,below=of engine] (catalogue) {
 \textbf{DemoCatalog}\\\rule{5.3cm}{0.3pt}\\
 product/photo database; encoder signature\\
 \rule{5.3cm}{0.3pt}\\\texttt{add}, \texttt{save}, \texttt{delete}\\
 \texttt{search\_with\_routing}};
\node[reportbox,text width=5.6cm,align=left,below=of segmenter] (backups) {
 \textbf{CatalogBackups}\\\rule{5.3cm}{0.3pt}\\
 store reference; session-bound previews\\
 \rule{5.3cm}{0.3pt}\\\texttt{export}, \texttt{preview}, \texttt{restore}\\
 recovery archive listing};
\node[reportbox,text width=5.6cm,align=left,below=of catalogue] (auth) {
 \textbf{AdminAuth}\\\rule{5.3cm}{0.3pt}\\
 account/session database; hasher; clock\\
 \rule{5.3cm}{0.3pt}\\\texttt{setup}, \texttt{login}, \texttt{session}\\
 \texttt{logout}, \texttt{change\_password}};
\node[reportbox,text width=5.6cm,align=left,below=of backups] (observe) {
 \textbf{SearchObservability}\\\rule{5.3cm}{0.3pt}\\
 independent telemetry database\\
 \rule{5.3cm}{0.3pt}\\\texttt{record\_search}, \texttt{record\_feedback}\\
 summary and failure-case reporting};
\draw[reportarrow,dashed] (backups.west)--node[above,font=\scriptsize]{uses}(catalogue.east);
\node[reportnote,text width=12.3cm,below=8mm of $(auth.south)!0.5!(observe.south)$] {
 FastAPI route handlers compose these services; arrows indicate dependencies, not inheritance.
 The added-catalogue encoder callback uses SearchEngine's fixed image encoder.};
\end{tikzpicture}

\caption{Implemented service classes and principal responsibilities}
\label{fig:components}\end{figure}

The product router receives an encoding callback rather than constructing a
second SigLIP2 model. New product photographs are prepared and encoded before
the transaction that changes the stored product. Retained photographs keep
their existing vectors. This avoids unnecessary repeated inference during a
metadata-only edit.

\section{Data and Database Design}
\subsection{Research-Gallery Metadata and Index Mapping}
Each searchable original image has an active metadata record containing its
relative path, file-derived product identifier, display title, category key,
and dataset source. The index mapping preserves the same image order as the
vectors. Loading checks that vector count, dimension, path count, metadata,
preprocessing, and catalogue checksum agree. A relative-path design allows the
gallery to be moved with its matching resources without embedding a personal
absolute path in every catalogue result.

The gallery image identifier is an application key. Reviewed groups, when
configured through the launcher, supply a separate mapping from specific images
to a reviewed visual identity. They do not replace the original metadata with
invented SKU labels.

\subsection{Managed Product Schema}
Figure~\ref{fig:product-er} shows the persisted product and image relationship.
The foreign key links each image to its product. The application enforces the
1--5-photo rule and unique ordered retained identifiers. SQLite foreign-key
checking is enabled on managed-catalogue connections. Product deletion explicitly
removes its images and then the product within the transaction.

\begin{figure}[H]\centering\begin{tikzpicture}[node distance=13mm]
\node[storebox,align=left,text width=5.7cm] (products) {
 \textbf{products}\\\rule{5.5cm}{0.3pt}\\
 \texttt{id} : text, primary key\\
 \texttt{name} : text\\
 \texttt{category} : text\\
 \texttt{price} : optional decimal text\\
 \texttt{currency} : text\\
 \texttt{available} : integer boolean\\
 \texttt{revision} : integer};
\node[storebox,align=left,text width=5.7cm,right=of products] (images) {
 \textbf{images}\\\rule{5.5cm}{0.3pt}\\
 \texttt{id} : text, primary key\\
 \texttt{product\_id} : text, foreign key\\
 \texttt{photo} : JPEG BLOB\\
 \texttt{vector} : float32 BLOB\\
 \texttt{position} : integer\\
 \phantom{revision}\\
 \textit{Application rule: 1--5 per product}};
\draw[reportarrow] (products.east)--node[above,font=\small]{1:N}(images.west);
\node[storebox,text width=11.5cm,below=of $(products.south)!0.5!(images.south)$] (config) {
 \textbf{config}: encoder/preprocessing signature and vector dimension\\
 Independent singleton configuration; not a product ownership relation};
\end{tikzpicture}

\caption{Physical product/photo schema and independent encoder configuration}
\label{fig:product-er}\end{figure}

Product UUIDs identify managed products, while image UUIDs identify their
photographs. The vector BLOB contains 768 float32 values, or 3,072 bytes before
database overhead. The \texttt{position} field records photograph order and the
first image is the cover. Prices are decimal text, preserving two fractional
digits without relying on binary floating-point currency arithmetic. An
integer revision supports conflict detection.

\subsection{Account and Monitoring Stores}
The account database (Figure~\ref{fig:auth-schema}) contains \path{admin_account},
\path{admin_sessions}, and \path{admin_login_failures}. A singleton
account record stores the username and Argon2 password hash. Session records
store a hash of the opaque session token, a CSRF token, and expiry time. No
per-product owner field or multiple seller identities are implied by this schema.

\begin{figure}[H]\centering\begin{tikzpicture}[node distance=10mm]
\node[storebox,text width=5.7cm,align=left] (account) {
 \textbf{admin\_account}\\\rule{5.5cm}{0.3pt}\\
 \texttt{id} : integer PK; CHECK(id=1)\\
 \texttt{username} : text, not null\\
 \texttt{password\_hash} : text, not null};
\node[storebox,text width=5.7cm,align=left,right=of account] (sessions) {
 \textbf{admin\_sessions}\\\rule{5.5cm}{0.3pt}\\
 \texttt{token\_hash} : text, PK\\
 \texttt{csrf\_token} : text, not null\\
 \texttt{expires\_at} : real, not null};
\node[storebox,text width=11.7cm,align=left,below=of $(account.south)!0.5!(sessions.south)$] (failures) {
 \textbf{admin\_login\_failures}\\
 \texttt{created\_at} : real, not null; no declared primary key};
\node[reportnote,text width=11.9cm,below=6mm of failures] {
 All three tables are in the independent account database. Sessions serve the
 singleton account; the physical schema has no account foreign-key column.
 No product-owner or multi-store relation is implied.};
\end{tikzpicture}

\caption{Physical account/session schema, separate from the product database}
\label{fig:auth-schema}\end{figure}

The monitoring database (Figure~\ref{fig:monitoring-er}) records query hashes, search modes, category routing,
ranked identifiers and scores, latency, and feedback. Its purpose is to inspect
local usage and failure cases without keeping uploaded search photographs.
Feedback is an observation supplied by a user, not an automatically verified
ground-truth label. The added-product search API returns query routing metadata
but does not share the original-gallery monitoring/feedback pipeline; this is
a current scope boundary.

\begin{figure}[H]\centering\begin{tikzpicture}[node distance=12mm]
\node[storebox,text width=5.9cm,align=left] (events) {
 \textbf{search\_events}\\\rule{5.7cm}{0.3pt}\\
 \texttt{search\_id} : text, PK\\
 \texttt{created\_at\_utc} : text\\
 \texttt{query\_sha256} : text\\
 \texttt{search\_mode} : text\\
 \texttt{requested\_category} : text\\
 \texttt{predicted\_category} : text?\\
 \texttt{applied\_category} : text?\\
 \texttt{category\_confidence} : real\\
 \texttt{top\_k} : integer\\
 \texttt{result\_count} : integer\\
 \texttt{latency\_ms} : real\\
 \texttt{top\_product\_id} : text?\\
 \texttt{top\_similarity} : real?\\
 \texttt{result\_product\_ids\_json} : text\\
 \texttt{result\_scores\_json} : text\\
 \texttt{model\_name} : text\\
 \texttt{preprocessing\_version} : text};
\node[storebox,text width=5.9cm,align=left,right=of events] (feedback) {
 \textbf{search\_feedback}\\\rule{5.7cm}{0.3pt}\\
 \texttt{feedback\_id} : text, PK\\
 \texttt{search\_id} : text, FK, UNIQUE\\
 \texttt{created\_at\_utc} : text\\
 \texttt{verdict} : text\\
 \texttt{selected\_product\_id} : text?\\
 \texttt{corrected\_category} : text?\\
 \texttt{notes} : text?};
\draw[reportarrow] (events.east)--node[above,font=\scriptsize]{1:0..1}(feedback.west);
\node[reportnote,text width=12.3cm,below=7mm of events.south east] {
 FK: search\_feedback.search\_id references search\_events.search\_id,
 ON DELETE CASCADE. ``?'' marks nullable fields. Product identifiers are recorded
 values, not cross-database SQL foreign keys. No raw query image is stored.};
\end{tikzpicture}

\caption{Search-event and feedback ERD: at most one feedback record per event}
\label{fig:monitoring-er}\end{figure}

\section{Retrieval and Selection Algorithms}
\subsection{Original-Gallery Retrieval}
The image encoder produces $\mathbf{z}_q$ as in Equation~\ref{eq:cosine}. For a
category $c$, the original-gallery router uses the normalised mean vector
\begin{equation}
 \boldsymbol{\mu}_c=
 \frac{\sum_{i:y_i=c}\mathbf{z}_i}
 {\left\lVert\sum_{i:y_i=c}\mathbf{z}_i\right\rVert_2},\qquad
 \widehat c=\arg\max_c \mathbf{z}_q^{\mathsf T}\boldsymbol{\mu}_c.
\end{equation}
Automatic mode applies $\widehat c$; explicit mode applies the requested category;
unrestricted mode applies none. FAISS ranks the full index, then the original
search engine filters candidates until the requested count is collected. When
reviewed product grouping is enabled, the handler obtains the relevant ranked
list and retains one result for each reviewed group. Unknown images retain their
own identifiers.

\subsection{Managed-Product Retrieval}
For a managed product $p$ with photograph set $V_p$, the effective score is
\begin{equation}
 S(q,p)=\max_{i\in V_p}s(q,i).
 \label{eq:product-max}
\end{equation}
The implementation ranks available photographs and keeps the first occurrence
of each product, which implements this maximum rule. Repeated photographs do not
accumulate votes and a product occupies one result card.

Automatic routing for this catalogue is conservative. It considers available
product photographs only. If any product lacks a category, fewer than two
categories exist, the best per-category photograph similarity is below 0.55,
or the margin between the best two category scores is below 0.08, it searches
all available products and returns a fallback reason. Otherwise it filters
to the winning category. These gates are implemented heuristics, not calibrated
probabilities or independently optimised thresholds.

\subsection{SAM2 Mask Handling}
The segmenter supplies all current points and binary labels to the local SAM2
processor. The model returns multiple candidate masks; the implementation
selects the highest predicted-quality candidate and restores it to the original
image size. Empty masks produce an error. The selected foreground is copied from
the original image into a white canvas bounded by the mask, and an RGBA green
overlay is returned separately for display.

The frontend preserves a usable prior selection if a later refinement fails,
but prevents a stale response from replacing a newer image or reset state.
The user's acceptance controls which cutout becomes the search query. The
quality score helps describe the mask proposal; the application has no annotated
mask reference against which to call it measured IoU accuracy.

\section{Integrity, Security, and Recovery Design}
\subsection{Administration Boundary}
A middleware guard normalises request paths and requires a valid session for
product management, backup, monitoring summaries, account changes, and reindexing.
Unsafe protected requests require a session-bound CSRF token and an acceptable
origin. Cookies are HttpOnly and SameSite Strict, and use the Secure flag when
served over HTTPS. The local HTTP deployment therefore does not claim encrypted
transport merely because these cookie settings exist.

Passwords are hashed through \texttt{argon2-cffi}; opaque random session tokens
are revocable and are represented in the database by SHA-256 hashes. Changing
the password invalidates existing sessions. First-account setup is limited to a
loopback client. The design is for a local singleton administrator and is not
a multi-tenant authorisation model.

\subsection{Product Transactions and Revisions}
Saving product details and photograph changes is one transaction. The server
checks that the expected revision matches, validates every retained/new image
reference, and increments the revision on success. A competing edit produces
HTTP 409 instead of silently replacing newer data. The photo preparation stage
occurs before the final commit, and the revision is checked again when the write
transaction begins.

\subsection{Backup and Restore}
A backup is a structured ZIP of JSON, JPEG photos, and float32 vectors, with
schema information, encoder compatibility data, and content hashes. The server
validates members, lengths, relationships, image order, image decoding, vector
dimension, finite values, and limits. It does not accept an uploaded SQLite
database or extract arbitrary ZIP paths to the filesystem. The maximum uploaded
archive is 200 MiB.

Preview creates a session-bound token valid for ten minutes and records the
current catalogue fingerprint. Confirmation is rejected if the catalogue has
changed. Before replacement, a recovery archive of the current catalogue is
written to disk. Restore then executes inside one SQLite transaction. A failure
to validate, write recovery data, or commit leaves the existing collection
intact. Product revisions advance so that editors opened before restore cannot
silently overwrite the restored state. Only the added-product catalogue is
restored; account, research-gallery, and telemetry data remain separate.

\section{Interface Design}
The main page brings together upload, source selection, search/selection mode,
category controls, result count, query preview, and ranked cards. The selection
tools are available for both catalogue sources. Original-gallery result cards
can provide relevance feedback, while managed-product cards carry the product's
stored metadata and available photographs.

The management page supports product entry, filtering, editing, image order and
cover choice, availability, and backup operations. Destructive restore has an
explicit preview and confirmation. Account settings are separate from product
details. Actual interface captures are presented with implementation evidence
in Chapter~\ref{ch:evaluation}; the cover artwork is a conceptual illustration,
not a screenshot or experiment result.

\section{Summary}
The design maps the required workflow onto separate selection, encoding,
ranking, product, account, and recovery components. Its data separation protects
the original research gallery, while transactions and revisions support a
maintainable local catalogue. The next chapter describes implementation and
the evidence obtained through testing and retrieval comparisons.
